\setcounterpageref{ex}{0}
\TaoPhieuKiemTra{20}

\setcounter{ex}{0}
\setcounter{bt}{0}

\textbf{\large I. PHẦN TRẮC NGHIỆM (3,0 điểm)}

{\itshape Khoanh tròn vào chữ cái đặt trước câu trả lời đúng. Mỗi câu đúng được 0,5 điểm.}

\begin{ex} % Câu 1
Số liền trước của số $100$ là
\choice
{$98$}
{$101$}
{$90$}
{\True $99$}
\loigiai{
Số liền trước của $100$ là $99$.\\
\textbf{Đáp án đúng: D}
}
\end{ex}

\begin{ex} % Câu 2
Biết số bị trừ là $88$, hiệu là $30$. Số trừ bằng
\choice
{$48$}
{\True $58$}
{$68$}
{$118$}
\loigiai{
Số trừ = Số bị trừ - Hiệu = $88 - 30 = 58$.\\
\textbf{Đáp án đúng: B}
}
\end{ex}

\begin{ex} % Câu 3
Chai nước ngọt chứa $2\text{ }l$, rót ra hai cốc mỗi cốc $1\text{ }l$ thì trong chai còn lại
\choice
{$1\text{ }l$}
{$2\text{ }l$}
{\True $0\text{ }l$ (hết nước)}
{$3\text{ }l$}
\loigiai{
Rót ra cả 2 cốc hết $1 + 1 = 2\text{ }l$, trong chai còn $2 - 2 = 0\text{ }l$.\\
\textbf{Đáp án đúng: C}
}
\end{ex}

\begin{ex} % Câu 4
Độ dài đường gấp khúc $ABC$ có $AB = 6\text{ cm}, BC = 8\text{ cm}$ là
\choice
{\True $14\text{ cm}$}
{$13\text{ cm}$}
{$15\text{ cm}$}
{$2\text{ cm}$}
\loigiai{
Độ dài: $6 + 8 = 14\text{ cm}$.\\
\textbf{Đáp án đúng: A}
}
\end{ex}

\begin{ex} % Câu 5
Điền dấu thích hợp: $52 + 16 \dots 70 - 2$.
\choice
{$>$}
{$<$}
{\True $=$}
{Không so sánh được}
\loigiai{
$52 + 16 = 68$; $70 - 2 = 68$. Vì $68 = 68$ nên điền dấu $=$.\\
\textbf{Đáp án đúng: C}
}
\end{ex}

\begin{ex} % Câu 6
Tính: $12 - 7 + 9 = \dots$
\choice
{$13$}
{$15$}
{$12$}
{\True $14$}
\loigiai{
$12 - 7 = 5$; $5 + 9 = 14$.\\
\textbf{Đáp án đúng: D}
}
\end{ex}

\vspace{4mm}
\textbf{\large II. PHẦN TỰ LUẬN (7,0 điểm)}

{\itshape Thí sinh trình bày chi tiết bài làm vào giấy kiểm tra.}

\begin{bt}[1,0 điểm]
Tính nhẩm:
\begin{multicols}{2}
\begin{enumerate}[a)]
  \item $8 + 4$
  \item $7 + 5$
  \item $12 - 5$
  \item $14 - 6$
\end{enumerate}
\end{multicols}
\loigiai{
\begin{itemize}
  \item $8 + 4 = 12$
  \item $7 + 5 = 12$
  \item $12 - 5 = 7$
  \item $14 - 6 = 8$
\end{itemize}
}
\end{bt}

\begin{bt}[1,5 điểm]
Đặt tính rồi tính:
\begin{multicols}{3}
\begin{enumerate}[a)]
  \item $33 + 42$
  \item $75 - 21$
  \item $56 + 56$
\end{enumerate}
\end{multicols}
\loigiai{
\begin{itemize}
  \item $33 + 42 = 75$
  \item $75 - 21 = 54$
  \item $56 + 56 = 112$
\end{itemize}
}
\end{bt}

\begin{bt}[1,0 điểm]
Tính giá trị biểu thức kèm đơn vị đo:
\begin{enumerate}[a)]
  \item $36\text{ kg} + 7\text{ kg} - 11\text{ kg}$
  \item $45\text{ }l - 11\text{ }l + 6\text{ }l$
\end{enumerate}
\loigiai{
\begin{enumerate}[a)]
  \item $36\text{ kg} + 7\text{ kg} - 11\text{ kg} = 43\text{ kg} - 11\text{ kg} = 32\text{ kg}$.
  \item $45\text{ }l - 11\text{ }l + 6\text{ }l = 34\text{ }l + 6\text{ }l = 40\text{ }l$.
\end{enumerate}
}
\end{bt}

\begin{bt}[1,5 điểm]
Bạn Vy có $35$ cây bút chì. Bạn Bảo có ít hơn bạn Vy $6$ cây bút chì. Hỏi bạn Bảo có bao nhiêu cây bút chì?
\loigiai{
Số cây bút chì bạn Bảo có là:
\[ 35 - 6 = 29 \text{ (cây bút chì)} \]
\textbf{Đáp số:} $29$ cây bút chì.
}
\end{bt}

\begin{bt}[1,0 điểm]
Một thùng chứa đầy $55\text{ }l$ dầu ăn. Buổi sáng cửa hàng đã bán ra $16\text{ }l$. Hỏi trong thùng còn lại bao nhiêu lít dầu ăn?
\loigiai{
Trong thùng còn lại số lít dầu ăn là:
\[ 55 - 16 = 39\text{ }l \]
\textbf{Đáp số:} $39\text{ }l$ dầu ăn.
}
\end{bt}

\begin{bt}[1,0 điểm]
Hình vẽ có đoạn thẳng $AB$. Nếu lấy thêm hai điểm $C$ và $D$ nằm giữa $A$ và $B$ thì có tất cả bao nhiêu đoạn thẳng?
\loigiai{
Các điểm thẳng hàng theo thứ tự là $A, C, D, B$.\\
Các đoạn thẳng tạo thành gồm:
- Đoạn đơn: $AC, CD, DB$ (3 đoạn).
- Đoạn đôi: $AD, CB$ (2 đoạn).
- Đoạn ba: $AB$ (1 đoạn).\\
Tổng cộng: $3 + 2 + 1 = 6$ đoạn thẳng.\\
\textbf{Đáp số:} $6$ đoạn thẳng.
}
\end{bt}
